\section{Related Work and Scope}\label{sec:related}
\subsection{Sensitivity-guided fault injection}
FKeras ranks edge-network parameters using gradients and leading weight-Hessian eigenpairs and orders bits by significance tiers~\cite{weng2024fkeras}. Our audit distinguishes its printed score from the effective sum strategy in the pinned public implementation. We preserve both definitions, adapt them to a continuous companion of the frozen KWS model, and keep the original tier-first ordering. This is source-level and numerical correspondence, not reproduction of the paper's original hardware or datasets.

PrioriFI uses Hessian-based initialization and observed fault outcomes to prioritize subsequent injections~\cite{weng2026priorifi}. Finding sensitive locations efficiently differs from minimizing expected risk of a fixed protected set under an exogenous fault law. Nevertheless, this is direct prior art on feedback-guided fault analysis and exceptions to significance monotonicity. Neither multiple-fault evaluation nor adaptive use of previous outcomes is novel here. We make no efficiency comparison against an unexecuted PrioriFI implementation.

\subsection{Protection mechanisms and learned allocations}
Catal\'an \emph{et al.} investigate fixed and variable bit protection using redundancy in stored numerical representations~\cite{catalan2025redundancy}. Those mechanisms and their hardware costs differ from forcing selected fault coordinates clean. COLA studies correlated feature errors and combines error coding with architectural/training changes~\cite{yu2023cola}. Its error-correlation object is not the anchored finite-bit interaction of a fixed executable.

The public \emph{Learned or Ranked?} companion release~\cite{alobaidi2026selectiveeccdataset} is particularly relevant. It provides learned and rule-based policies, cost-matched held-out comparisons, BER sweeps, and layer-interaction diagnostics. The source snapshot includes layer/subset additivity and contextual marginal analyses. This rules out a broad claim that protection dependence or layer interactions have not been studied. Its layer-wise choices, word-level ECC variants, burst/interleaving laws, and accuracy--cost reward differ from our fixed-cardinality encoded-bit problem. We do not assert that its layer interactions equal our candidate-bit M\"obius coefficients, or that ideal suppression beats ECC. The release is a dataset; peer-reviewed publication of the accompanying manuscript has not been established in our audit.

\subsection{Established mathematical foundations}
Multilinear extensions and biased Fourier analysis of functions on the Boolean cube are established tools~\cite{odonnell2014analysis}. Classical reliability importance also uses the difference between forced component states averaged over the remaining system~\cite{rauzy2022reliability}. Our signed, real-valued task risk need not be monotone, unlike a coherent failure indicator. Applying these identities to encoded neural-network faults does not make the identities new. Our claimed contribution is their operational separation and controlled validation on exact executables, including selection versus independent evaluation, nonadditive task metrics, and explicit approximation limits.

The resulting comparison is deliberately not a claim that static methods have no interactions and conditional methods do. A continuous Hessian contains cross-parameter curvature. A finite single-bit jump can be large; a projected parameter score need not estimate a probability-space derivative conditioned on protection. Conversely, an exact conditional measurement on a small candidate pool does not establish a globally optimal protection policy. These distinctions determine the baselines and claim boundaries throughout the paper.
